\chapter{Sucesiones y series de funciones}

El estudio de sucesiones y series de funciones extiende el análisis numérico al espacio
funcional. A diferencia de las series numéricas, aquí la convergencia depende del punto
$x$ en el dominio, lo que introduce dos nociones fundamentalmente distintas: convergencia
\emph{puntual} y convergencia \emph{uniforme}. Esta distinción es crucial, pues solo la
convergencia uniforme garantiza que propiedades analíticas esenciales ---continuidad,
integrabilidad, derivabilidad--- se preserven al pasar al límite.

% ------------------------------------------------------------
\section{Sucesiones de funciones}

\begin{definition}[Sucesión de funciones]
Sea $X \subset \mathbb{R}$. Una \textbf{sucesión de funciones} $\{f_n\}_{n=1}^\infty$ en
$X$ es una familia de funciones $f_n : X \to \mathbb{R}$ indexada por $n \in \mathbb{N}$.
\end{definition}

\begin{definition}[Convergencia puntual y uniforme]
Sea $\{f_n\}$ una sucesión de funciones en $X$ y $f : X \to \mathbb{R}$.
\begin{enumerate}
    \item $\{f_n\}$ \textbf{converge puntualmente} a $f$ en $X$ si para cada
          $x_0 \in X$ fijo se tiene $\lim_{n \to \infty} f_n(x_0) = f(x_0)$, es decir:
          \[
              \forall x \in X,\;\forall\varepsilon > 0,\;
              \exists N(\varepsilon,x) \in \mathbb{N}:\;
              n > N \implies |f_n(x) - f(x)| < \varepsilon.
          \]
    \item $\{f_n\}$ \textbf{converge uniformemente} a $f$ en $X$ si la convergencia es
          simultánea en todo $X$:
          \[
              \forall\varepsilon > 0,\;
              \exists N(\varepsilon) \in \mathbb{N}:\;
              n > N \implies \sup_{x\in X}|f_n(x) - f(x)| < \varepsilon.
          \]
\end{enumerate}
\end{definition}

\begin{rem}
La diferencia técnica clave reside en el \emph{orden de los cuantificadores}:
\begin{align*}
    \text{Puntual:}  &\quad \forall x\;\forall\varepsilon\;\exists N(x,\varepsilon)\colon\ldots\\
    \text{Uniforme:} &\quad \forall\varepsilon\;\exists N(\varepsilon)\;\forall x\colon\ldots
\end{align*}
En la convergencia uniforme, \textbf{un solo $N$ sirve para todos los $x \in X$
simultáneamente}, lo que tiene consecuencias analíticas profundas.
\end{rem}

\begin{example}[Convergencia puntual sin uniforme: $f_n(x) = x^n$]
En $X = [0,1]$, la función límite de $f_n(x) = x^n$ es
\[
    f(x) = \begin{cases} 0 & 0 \leq x < 1, \\ 1 & x = 1. \end{cases}
\]
Cada $f_n$ es continua, pero $f$ es discontinua en $x = 1$. Por el
Teorema~\ref{thm:prop_uniforme}, la convergencia no puede ser uniforme. Directamente:
\[
    \sup_{x\in[0,1]} |x^n - f(x)| = \sup_{x\in[0,1)} x^n = 1 \not\to 0.
\]
Sin embargo, para $0 < r < 1$ fijo, $\sup_{x\in[0,r]} x^n = r^n \to 0$, de modo que
la convergencia \textbf{sí} es uniforme en $[0, r]$.
\end{example}

\begin{example}[Convergencia uniforme]
Sea $f_n(x) = x/n$ en $[0,1]$. Entonces $f(x) = 0$ y
$\sup_{x\in[0,1]}|x/n| = 1/n \to 0$. La convergencia es uniforme en $[0,1]$.
\end{example}

\begin{example}[Desplazamiento del máximo]
Sea $f_n(x) = nxe^{-nx}$ en $[0,+\infty)$. Para cada $x > 0$ fijo, $f_n(x) \to 0$;
y $f_n(0) = 0$, luego la convergencia es puntual a $0$. Sin embargo, el máximo se
alcanza en $x = 1/n$ con $f_n(1/n) = e^{-1}$, de modo que
$\sup_{x \geq 0} |f_n(x)| = e^{-1} \not\to 0$: la convergencia
\textbf{no} es uniforme en $[0, +\infty)$.
\end{example}

\begin{theorem}[Criterio de Cauchy para convergencia uniforme]
$\{f_n\}$ converge uniformemente en $X$ si y solo si para todo $\varepsilon > 0$ existe
$N \in \mathbb{N}$ tal que $|f_n(x) - f_m(x)| < \varepsilon$ para todo $n, m > N$ y
todo $x \in X$.
\end{theorem}

\begin{theorem}[Propiedades de la convergencia uniforme]
\label{thm:prop_uniforme}
Sea $\{f_n\}$ una sucesión que converge uniformemente a $f$ en $[a,b]$.
\begin{enumerate}
    \item \textbf{Continuidad.} Si cada $f_n$ es continua en $[a,b]$, entonces $f$ es
          continua en $[a,b]$. En particular,
          $\lim_{n\to\infty}\lim_{x\to x_0} f_n(x)
          = \lim_{x\to x_0}\lim_{n\to\infty} f_n(x)$.
    \item \textbf{Integrabilidad.} Si cada $f_n$ es Riemann-integrable en $[a,b]$,
          entonces $f$ es integrable y
          \[
              \lim_{n \to \infty} \int_a^b f_n(x)\,dx
              = \int_a^b \lim_{n\to\infty} f_n(x)\,dx
              = \int_a^b f(x)\,dx.
          \]
    \item \textbf{Derivabilidad.} Si cada $f_n \in C^1[a,b]$, $\{f_n(x_0)\}$
          converge para algún $x_0 \in [a,b]$ y $\{f_n'\}$ converge uniformemente
          en $[a,b]$, entonces $f$ es derivable y
          $f'(x) = \lim_{n \to \infty} f_n'(x)$ para todo $x \in [a,b]$.
\end{enumerate}
\end{theorem}

\begin{rem}
Estos resultados pueden resumirse así: \emph{la convergencia uniforme conmuta con los
límites, integrales y derivadas}. La mera convergencia puntual no garantiza ninguna de
estas propiedades. El ejemplo $f_n(x) = x^n$ en $[0,1]$ ilustra la pérdida de
continuidad; $g_n(x) = \sin(nx)/\sqrt{n}$ ilustra la pérdida de derivabilidad en el
límite.
\end{rem}

\begin{example}[Intercambio inválido de derivada y límite]
Sea $f_n(x) = \dfrac{\sin(nx)}{n}$ en $\mathbb{R}$. Entonces $f_n \to 0$
uniformemente (pues $|f_n(x)| \leq 1/n$), con $f \equiv 0$. Sin embargo,
$f_n'(x) = \cos(nx)$, que no converge en ningún $x$ genérico. El Teorema
\ref{thm:prop_uniforme}.3 no se aplica porque $\{f_n'\}$ no converge uniformemente.
\end{example}

\begin{theorem}[Teorema de Dini]
Sea $K \subset \mathbb{R}$ compacto y $\{f_n\}$ una sucesión de funciones continuas en
$K$ que converge puntualmente a una función continua $f$. Si la sucesión es monótona
(es decir, $f_n \leq f_{n+1}$ en $K$ para todo $n$, o al revés), entonces la
convergencia es uniforme en $K$.
\end{theorem}

\begin{rem}
Las hipótesis del teorema de Dini son todas necesarias: existen contraejemplos si se
omite la compacidad, la monotonía o la continuidad del límite.
\end{rem}

\begin{theorem}[Teorema de aproximación de Weierstrass]
Sea $f \in C[a,b]$. Para todo $\varepsilon > 0$ existe un polinomio $P(x)$ tal que
$\sup_{x\in[a,b]}|f(x) - P(x)| < \varepsilon$.
\end{theorem}

\begin{rem}
El teorema garantiza que los polinomios son densos en $C[a,b]$ con la norma del supremo,
lo que fundamenta métodos de aproximación como la interpolación de Lagrange, los
polinomios de Bernstein y los splines.
\end{rem}

% ------------------------------------------------------------
\section{Series de funciones}

\begin{definition}[Serie de funciones]
Dada $\{u_n(x)\}$ en $X$, la \textbf{serie de funciones}
$\displaystyle\sum_{n=1}^\infty u_n(x)$ se define por sus sumas parciales
$S_N(x) = \displaystyle\sum_{n=1}^N u_n(x)$. La serie converge puntual (resp.
uniformemente) a $S(x)$ si $\{S_N(x)\}$ converge puntual (resp. uniformemente) a
$S(x)$.
\end{definition}

\begin{theorem}[Prueba M de Weierstrass]
Si existe una sucesión numérica $\{M_n\}$ tal que $|u_n(x)| \leq M_n$ para todo
$x \in X$ y $\displaystyle\sum_{n=1}^\infty M_n < \infty$, entonces
$\displaystyle\sum_{n=1}^\infty u_n(x)$ converge \textbf{absoluta y uniformemente}
en $X$.
\end{theorem}

\begin{example}
La serie $\displaystyle\sum_{n=1}^\infty \frac{\cos(nx)}{n^2}$ en $\mathbb{R}$:
$\left|\dfrac{\cos(nx)}{n^2}\right| \leq \dfrac{1}{n^2}$ y $\sum 1/n^2$ converge.
Por la prueba M, la serie converge absoluta y uniformemente en $\mathbb{R}$.
En particular, su suma es una función continua.
\end{example}

\begin{example}
La serie $\displaystyle\sum_{n=1}^\infty \frac{x^n}{n^2}$ en $[-1,1]$:
$\left|\dfrac{x^n}{n^2}\right| \leq \dfrac{1}{n^2}$ para $|x| \leq 1$.
La prueba M da convergencia uniforme en $[-1,1]$ y permite derivar e integrar
término a término en $(-1,1)$.
\end{example}

\begin{rem}
La prueba M es suficiente pero no necesaria para convergencia uniforme. Una serie puede
converger uniformemente sin admitir una mayorante $\sum M_n$ convergente.
\end{rem}

\begin{theorem}[Criterios de Dirichlet y Abel para convergencia uniforme]
Sean $\{a_n(x)\}$ y $\{b_n(x)\}$ funciones en $X$.
\begin{itemize}
    \item \textbf{Dirichlet:} Si las sumas parciales de $\sum b_n(x)$ están
          uniformemente acotadas en $X$ y $\{a_n(x)\}$ tiende a $0$ uniformemente
          y es monótona en $n$ para cada $x$, entonces
          $\sum a_n(x)b_n(x)$ converge uniformemente en $X$.
    \item \textbf{Abel:} Si $\sum b_n(x)$ converge uniformemente y
          $\{a_n(x)\}$ es uniformemente acotada y monótona en $n$ para cada $x$,
          entonces $\sum a_n(x)b_n(x)$ converge uniformemente.
\end{itemize}
\end{theorem}

\begin{theorem}[Integración y derivación término a término]
Si $\sum u_n(x)$ converge uniformemente en $[a,b]$ y cada $u_n$ es continua:
\[
    \int_a^b \sum_{n=1}^\infty u_n(x)\,dx
    = \sum_{n=1}^\infty \int_a^b u_n(x)\,dx.
\]
Si además $\sum u_n'(x)$ converge uniformemente en $[a,b]$ y $\sum u_n(x_0)$ converge
para algún $x_0$:
\[
    \frac{d}{dx}\sum_{n=1}^\infty u_n(x) = \sum_{n=1}^\infty u_n'(x).
\]
\end{theorem}

\begin{example}[La función de Weierstrass]
La función $W(x) = \displaystyle\sum_{n=0}^\infty a^n\cos(b^n\pi x)$, con $0 < a < 1$,
$b$ entero impar y $ab > 1 + \dfrac{3\pi}{2}$, es continua en $\mathbb{R}$ (prueba M:
$|a^n\cos(\cdot)| \leq a^n$ y $\sum a^n < \infty$) pero \textbf{no es derivable en
ningún punto}. Publicada por Weierstrass en 1872, demolió la creencia de que toda
función continua es derivable salvo en un conjunto excepcional.
\end{example}

% ------------------------------------------------------------
\section{Series de potencias}

\begin{definition}[Serie de potencias]
Una \textbf{serie de potencias} centrada en $a \in \mathbb{R}$ tiene la forma:
\[
    \sum_{n=0}^\infty c_n(x-a)^n
    = c_0 + c_1(x-a) + c_2(x-a)^2 + \cdots
\]
donde $\{c_n\}$ son constantes reales.
\end{definition}

\begin{theorem}[Radio e intervalo de convergencia]
Para toda serie de potencias existe un único $R \in [0, +\infty]$ (el
\textbf{radio de convergencia}) tal que:
\begin{itemize}
    \item La serie converge absolutamente si $|x - a| < R$.
    \item La serie diverge si $|x - a| > R$.
    \item En $|x - a| = R$ el comportamiento debe analizarse por separado.
\end{itemize}
El radio se obtiene mediante la \textbf{fórmula de Cauchy-Hadamard}:
\[
    R = \frac{1}{\displaystyle\limsup_{n \to \infty} \sqrt[n]{|c_n|}},
\]
o bien, si el límite existe, $R = \displaystyle\lim_{n \to \infty}
\left|\dfrac{c_n}{c_{n+1}}\right|$.
\end{theorem}

\begin{example}[Cálculo del radio de convergencia]
\begin{enumerate}
    \item $\displaystyle\sum_{n=0}^\infty \frac{x^n}{n!}$:\;
          $R = \lim \dfrac{1/n!}{1/(n+1)!} = \lim(n+1) = +\infty$.
          Converge para todo $x \in \mathbb{R}$.
    \item $\displaystyle\sum_{n=0}^\infty n!\,x^n$:\;
          $R = \lim \dfrac{n!}{(n+1)!} = 0$.
          Converge solo en $x = 0$.
    \item $\displaystyle\sum_{n=1}^\infty \frac{x^n}{n}$:\;
          $R = 1$. En $x = 1$: serie armónica (diverge). En $x = -1$: serie alternada
          armónica (converge). Intervalo: $[-1, 1)$.
    \item $\displaystyle\sum_{n=0}^\infty (-1)^n\frac{x^{2n}}{(2n)!}$:\;
          $R = +\infty$ (serie del coseno).
\end{enumerate}
\end{example}

\begin{theorem}[Convergencia uniforme en compactos]
Si $\sum c_n(x-a)^n$ tiene radio $R > 0$, entonces converge \textbf{absoluta y
uniformemente} en todo intervalo cerrado $[a-r, a+r]$ con $0 < r < R$.
\end{theorem}

\begin{rem}
La convergencia uniforme no está garantizada en el intervalo abierto completo $(-R+a,\,
R+a)$. Por ejemplo, $\sum x^n/n$ tiene $R = 1$ y converge uniformemente en $[-r, r]$
para cada $r < 1$, pero no uniformemente en $(-1, 1)$. En los extremos $x = a \pm R$
la convergencia debe verificarse por separado.
\end{rem}

\begin{theorem}[Analiticidad: derivación e integración término a término]
Dentro del intervalo $|x - a| < R$, la suma $f(x) = \sum_{n=0}^\infty c_n(x-a)^n$ es
infinitamente diferenciable y:
\[
    f'(x) = \sum_{n=1}^\infty n\,c_n(x-a)^{n-1},
    \qquad
    \int_a^x f(t)\,dt = \sum_{n=0}^\infty \frac{c_n}{n+1}(x-a)^{n+1}.
\]
Ambas series tienen el mismo radio de convergencia $R$.
\end{theorem}

\begin{theorem}[Serie de Taylor]
Si $f$ tiene representación en serie de potencias en un entorno de $a$, dicha
representación es única y sus coeficientes son $c_n = \dfrac{f^{(n)}(a)}{n!}$. La
\textbf{serie de Taylor} de $f$ en $a$ es:
\[
    \sum_{n=0}^\infty \frac{f^{(n)}(a)}{n!}(x-a)^n.
\]
La serie converge a $f(x)$ si y solo si el resto de Lagrange
$R_n(x) = \dfrac{f^{(n+1)}(\xi)}{(n+1)!}(x-a)^{n+1}$ satisface $R_n(x) \to 0$
cuando $n \to \infty$ (para algún $\xi$ entre $a$ y $x$).
\end{theorem}

\begin{example}[Expansiones estándar]
\begin{itemize}
    \item $e^x = \displaystyle\sum_{n=0}^\infty \frac{x^n}{n!},\quad x \in \mathbb{R}$.
    \item $\sin x = \displaystyle\sum_{n=0}^\infty \frac{(-1)^n x^{2n+1}}{(2n+1)!},
          \quad \cos x = \displaystyle\sum_{n=0}^\infty \frac{(-1)^n x^{2n}}{(2n)!},
          \quad x \in \mathbb{R}$.
    \item $\ln(1+x) = \displaystyle\sum_{n=1}^\infty \frac{(-1)^{n-1}x^n}{n},
          \quad x \in (-1,1]$.
    \item $\dfrac{1}{1-x} = \displaystyle\sum_{n=0}^\infty x^n, \quad |x| < 1$.
    \item $\arctan x = \displaystyle\sum_{n=0}^\infty \frac{(-1)^n x^{2n+1}}{2n+1},
          \quad |x| \leq 1$.
    \item $(1+x)^\alpha = \displaystyle\sum_{n=0}^\infty \binom{\alpha}{n}x^n,
          \quad |x| < 1$, \quad donde
          $\binom{\alpha}{n} = \dfrac{\alpha(\alpha-1)\cdots(\alpha-n+1)}{n!}$.
\end{itemize}
\end{example}

\begin{example}[Aplicaciones de las series de Taylor]
\begin{enumerate}
    \item \textbf{Límites indeterminados:}
    \[
        \lim_{x \to 0} \frac{\sin x - x}{x^3}
        = \lim_{x \to 0}
          \frac{\bigl(x - \tfrac{x^3}{6} + O(x^5)\bigr) - x}{x^3}
        = -\frac{1}{6}.
    \]
    \item \textbf{Indeterminación $1^\infty$:}
    \[
        \lim_{x \to 0}(\cos x)^{1/x^2}
        = \lim_{x \to 0} e^{\,\ln(\cos x)/x^2}.
    \]
    Como $\ln(\cos x) = \ln\!\bigl(1 - \tfrac{x^2}{2} + O(x^4)\bigr)
    \sim -\tfrac{x^2}{2}$, el exponente tiende a $-1/2$ y el límite es
    $e^{-1/2}$.
    \item \textbf{Integral sin primitiva elemental:}
    \[
        \int_0^1 e^{-x^2}\,dx
        = \sum_{n=0}^\infty \frac{(-1)^n}{n!\,(2n+1)}
        \approx 0{,}7468\ldots
    \]
    \item \textbf{Estimación del error:} Aproximar $e$ con error $< 10^{-4}$.
    El resto de la serie de $e^1$ satisface
    $R_N < \dfrac{3}{(N+1)!}$. Para $N = 8$:
    $3/9! \approx 8\times 10^{-6} < 10^{-5}$.
    \item \textbf{Ecuaciones diferenciales:} La ecuación de Airy $y'' = xy$
    se resuelve sustituyendo $y = \sum_{n=0}^\infty a_n x^n$ y determinando
    los coeficientes por recurrencia: $a_{n+3} = \dfrac{a_n}{(n+2)(n+3)}$,
    con condiciones iniciales $a_0 = y(0)$ y $a_1 = y'(0)$.
\end{enumerate}
\end{example}

% ------------------------------------------------------------
\section{Complementos avanzados}

\begin{theorem}[Irracionalidad de $e$]
El número $e = \sum_{n=0}^\infty 1/n!$ es irracional. Si $e = p/q$, entonces
\[
    0 < q!\!\left(e - \sum_{n=0}^q \frac{1}{n!}\right)
    = \sum_{k=1}^\infty \frac{1}{(q+1)(q+2)\cdots(q+k)}
    < \sum_{k=1}^\infty \frac{1}{(q+1)^k} = \frac{1}{q} \leq 1,
\]
contradiciendo que la expresión de la izquierda sería un entero positivo.
\end{theorem}

\begin{theorem}[Teorema de Borel]
Dada cualquier sucesión $\{a_n\}_{n=0}^\infty$ de números reales, existe una función
$f \in C^\infty(\mathbb{R})$ tal que $f^{(n)}(0) = a_n$ para todo $n$. Es decir,
cualquier sucesión puede ser la sucesión de coeficientes de Taylor de alguna función
infinitamente diferenciable.
\end{theorem}

\begin{rem}
El teorema de Borel muestra que la clase $C^\infty$ es mucho más amplia que la de las
funciones analíticas. Tener derivadas de todos los órdenes \textbf{no} implica que la
serie de Taylor converja a la función.
\end{rem}

\begin{example}[Función $C^\infty$ no analítica]
Defínase
\[
    \varphi(x) = \begin{cases} e^{-1/x^2} & x \neq 0, \\ 0 & x = 0. \end{cases}
\]
Por inducción se demuestra que $\varphi^{(n)}(0) = 0$ para todo $n \in \mathbb{N}_0$.
Luego su serie de Taylor en $0$ es la serie nula, que converge a $0$ para todo $x$,
pero $\varphi(x) \neq 0$ para $x \neq 0$. Así, $\varphi \in C^\infty(\mathbb{R})$
pero \textbf{no es analítica en $0$}.

Esta función es la piedra de toque con la que se construyen las funciones meseta y las
particiones de la unidad, herramientas indispensables en análisis diferencial y teoría
de distribuciones.
\end{example}

\begin{rem}
La condición $R_n(x) \to 0$ (convergencia del resto de Taylor) es una hipótesis
adicional, no automática para funciones de clase $C^\infty$. La
\textbf{analiticidad} es una propiedad más restrictiva y ``rígida'' que la
diferenciabilidad infinita.
\end{rem}

\begin{theorem}[Unicidad y álgebra de series de potencias]
Si $\sum a_n(x-c)^n = \sum b_n(x-c)^n$ en algún intervalo con centro $c$, entonces
$a_n = b_n$ para todo $n$. Dentro del radio de convergencia, las series de potencias
pueden:
\begin{itemize}
    \item \textbf{Sumarse:} con radio $\geq \min(R_1,R_2)$.
    \item \textbf{Multiplicarse} (producto de Cauchy): con radio $\geq \min(R_1,R_2)$.
    \item \textbf{Componerse:} si $g(0) = 0$ (o más generalmente $|g(x)| < R_f$),
          entonces $f(g(x))$ se puede expandir en serie de potencias con radio
          positivo.
\end{itemize}
\end{theorem}

\begin{rem}
La teoría de series de potencias cierra un ciclo conceptual que une la aritmética de los
reales, el análisis de funciones y los métodos de aproximación. Estas herramientas son
la base del análisis complejo, la teoría de ecuaciones diferenciales ordinarias, la
física matemática y los métodos numéricos modernos.
\end{rem}